\chapter{Exemplos Práticos}\label{cap:exemplos}

Este capítulo reúne três cenários que combinam as ferramentas dos capítulos anteriores. Cada um parte de uma pergunta de diagnóstico e usa só comandos já apresentados, de modo que o leitor possa reproduzi-los em um terminal com dois painéis: um para a ferramenta de coleta e outro para gerar a atividade observada. Todos exigem \texttt{root} e foram testados no ambiente descrito no capítulo \ref{cap:monitoramento}.

\section{Monitoramento de processos}

\textbf{Pergunta:} que processos de vida curta estão sendo criados no sistema? Um \textit{script} que roda por milissegundos, disparado pelo \texttt{cron} ou por um serviço, não aparece em um \texttt{ps} nem em uma amostra de \textit{polling}.

No primeiro terminal, inicie o rastreamento de \texttt{execve} e deixe-o rodando:

\begin{lstlisting}
sudo execsnoop-libbpf
\end{lstlisting}

No segundo terminal, crie alguns processos curtos:

\begin{lstlisting}
for i in 1 2 3; do date > /dev/null; sleep 0.2; done
\end{lstlisting}

O primeiro terminal passa a listar cada \texttt{date} e cada \texttt{sleep}, com o PID, o processo pai e os argumentos. O \texttt{ps} não mostraria nenhum deles. Para saber \textit{quem} está criando processos em maior volume, em vez de listar cada um, troque a ferramenta por uma contagem:

\begin{lstlisting}
sudo bpftrace -e 'tracepoint:syscalls:sys_enter_execve { @[comm] = count(); }'
\end{lstlisting}

Ao encerrar com \texttt{Ctrl+C}, o \texttt{bpftrace} imprime o mapa com o total de execuções por processo pai. Um número alto em um \textit{script} de verificação periódica é um indício de que ele cria processos demais para o que faz.

\section{Análise de chamadas de sistema}

\textbf{Pergunta:} o que um programa está pedindo ao kernel? Aqui o alvo é o servidor HTTP embutido do Python. Inicie-o no primeiro terminal:

\begin{lstlisting}
python3 -m http.server 8080
\end{lstlisting}

No segundo terminal, descubra o PID do servidor e observe os arquivos que ele abre enquanto atende requisições:

\begin{lstlisting}
PID=$(pgrep -f "http.server 8080")
sudo opensnoop-libbpf -p $PID
\end{lstlisting}

Em um terceiro terminal, faça requisições com \texttt{curl http://localhost:8080/}. Cada requisição aparece no \texttt{opensnoop} como a abertura do arquivo ou do diretório pedido, e é possível ver o que o servidor lê do disco para responder.

Para a visão agregada, o \texttt{syscount} mostra quais chamadas o servidor mais faz, sem interromper o processo, o que o \texttt{strace} não consegue evitar:

\begin{lstlisting}
sudo syscount-libbpf -p $PID -i 5
\end{lstlisting}

Comparar essa saída antes e depois de uma carga de requisições revela a proporção entre chamadas de rede (\texttt{sendto}, \texttt{recvfrom}), de arquivo (\texttt{openat}, \texttt{read}) e de espera (\texttt{epoll\_wait}, \texttt{poll}). Uma espera dominante indica um servidor ocioso, não lento.

\section{Rastreamento de uma aplicação}

\textbf{Pergunta:} quem se conecta ao serviço, por quanto tempo e com quantos bytes? Com o mesmo servidor na porta 8080, rode o \texttt{tcplife}, que registra uma linha ao fim de cada conexão:

\begin{lstlisting}
sudo tcplife-libbpf -L 8080
\end{lstlisting}

Gere conexões em outro terminal:

\begin{lstlisting}
for i in $(seq 1 5); do curl -s http://localhost:8080/ > /dev/null; done
\end{lstlisting}

Cada linha traz o processo, os endereços e portas, os bytes enviados e recebidos e a duração. Conexões longas ou com muitos bytes se destacam do conjunto. Para ver também quem \textit{abre} conexões a partir da máquina, e não só as recebidas, use o \texttt{tcpconnect-libbpf} em paralelo.

Esses comandos observam o serviço de fora, pelo que o kernel vê. Para olhar por dentro de um processo Python, sem alterá-lo, use a sonda USDT do interpretador (seção \ref{sec:uprobes} e capítulo \ref{cap:instrumentacao}). O comando a seguir lista os módulos importados por qualquer programa Python, e ajuda a explicar um início lento:

\begin{lstlisting}
sudo bpftrace -e 'usdt:/usr/bin/python3.14:python:import__find__load__start {
    printf("%s\n", str(arg0)); }'
\end{lstlisting}

Em outro terminal, execute \texttt{python3 -c "import json, csv"} e observe os módulos aparecerem na ordem em que o interpretador os carrega. Para transformar uma dessas medidas em uma série contínua em um painel, siga o caminho do capítulo \ref{cap:coleta}.
